\labsection{10}{Searching and sorting experiments}{sec:lab10}

\begin{sourcealignment}{Official Lab Manual --- Lab X}
\textbf{Objective.} Search a list and perform sorting using different techniques.

\textbf{Outcomes.} Understand searching techniques and the working of multiple sorting techniques.

\textbf{Manual problems.}
\begin{enumerate}
  \item Sort an array in descending order using quick sort, always selecting the median element of the current array/subarray as pivot.
  \item Implement binary search and test whether a given element is present.
  \item Modify merge logic to merge three sorted arrays into one array.
  \item Sort an array using heap sort.
\end{enumerate}
\end{sourcealignment}

\begin{labproject}{Trace, implement, and compare algorithms}
\textbf{Coverage.} Chapter 10.\\
\textbf{Goal.} Implement binary search and several sorting methods, then compare them on one common dataset.

Use a common dataset where appropriate, for example:
\begin{lstlisting}
vector<int> data = {29, 10, 14, 37, 13};
\end{lstlisting}

\begin{tryit}
Trace at least two sorting methods by hand. Then implement the four manual problems above. Count comparisons or swaps for one sorting pair so the final-project style of algorithm comparison becomes familiar.
\end{tryit}

\begin{clarification}
The source's quick-sort theory contains a complexity inconsistency, documented in Chapter~\ref{ch:search-sort}. Also, ``median element as pivot'' can mean the element at the middle index rather than the statistical median value; follow the wording used by your instructor or clarify the interpretation before coding.
\end{clarification}
\end{labproject}

\begin{verification}
\begin{itemize}
  \item Binary search is run only on sorted input.
  \item Descending quick sort produces a correctly ordered final array.
  \item The three-way merge preserves every input value exactly once.
  \item Heap sort output is sorted and the heap property is maintained during extraction.
  \item Test cases include duplicates and a missing search target.
\end{itemize}
\end{verification}

\begin{labdeliverable}
Submit implementations for the assigned manual problems, sample output, and at least one complete trace showing intermediate states.
\end{labdeliverable}
